% ═══════════════════════════════════════════════════════════════
%  chapters/ch15-lifecycle.tex
%  Жизненный цикл ПК-приложения
% ═══════════════════════════════════════════════════════════════

\chapter{Жизненный цикл приложения}
\label{ch:lifecycle}

\epigraph{<<Хорошая программа не только хорошо работает.
Она ещё и правильно запускается, и правильно завершается>>.}{}

\section{Введение}
\label{sec:lifecycle-intro}

В \texttt{ch12}--\texttt{ch14} мы разобрали \emph{потоки},
\emph{отрисовку}, \emph{интерфейс} \emph{генератора}. \emph{Но}
\emph{осталось} \emph{главное}~--- \emph{как} \emph{всё} \emph{это}
\emph{запустить} \emph{и} \emph{остановить} \emph{правильно}.

\textbf{Жизненный цикл}~--- \emph{это} \emph{порядок}. \emph{От}
\emph{первого} \emph{вызова} \texttt{main()} \emph{до}
\emph{последнего} \texttt{return 0}.

\begin{ideabox}
\textbf{Жизненный цикл~--- это \emph{порядок}.} \emph{Есть}
\emph{правильный} \emph{порядок} \emph{инициализации} \emph{и}
\emph{правильный} \emph{порядок} \emph{завершения}. \emph{Если}
\emph{нарушить}~--- \emph{будут} \emph{странные} \emph{баги}:
\emph{утечки}, \emph{зависания}, \emph{<<плавающие>>} \emph{ошибки}.
\end{ideabox}

\subsection{Что вы узнаете}

\begin{itemize}
    \item Что делает \texttt{main()}.
    \item В каком порядке инициализируются модули.
    \item Почему \texttt{app\_data\_free}~--- не <<просто
          освобождение памяти>>.
    \item Как правильно останавливать таймеры.
    \item Почему \texttt{delete-event} и \texttt{destroy}~---
          разные вещи.
    \item Что произойдёт, если \emph{нарушить} порядок.
\end{itemize}

\section{Точка входа: \texttt{main()}}
\label{sec:lifecycle-main}

\texttt{main.c}~--- это самый маленький файл проекта.
В нём всего 30 строк. Но каждая строка — важна.
Здесь задаётся порядок всей жизни приложения.

\begin{listing}[htbp]
\begin{minted}[fontsize=\footnotesize, frame=none, bgcolor=gray!5]{c}
int main(int argc, char *argv[])
{
    debug_log_init();
    LOG_DEBUG("=== ПРОГРАММА ЗАПУЩЕНА ===");

    gtk_init(&argc, &argv);

    AppData app;
    app_data_init(&app);
    load_settings(&app);

    ui_create(&app);
    scan_serial_ports(&app);

    gtk_widget_show_all(app.window);
    gtk_main();

    save_settings(&app);
    app_data_free(&app);

    LOG_DEBUG("=== ПРОГРАММА ЗАВЕРШЕНА ===");
    return 0;
}
\end{minted}
\caption{\texttt{main()}~--- точка входа}
\label{lst:lifecycle-main}
\end{listing}

\subsection{Порядок вызовов}

\textbf{Одиннадцать строк} — одиннадцать решений.
Порядок не случаен:

\begin{enumerate}
    \item \texttt{debug\_log\_init()}~--- до всего.
          Если \texttt{gtk\_init} упадёт~--- мы хотим
          знать почему. Значит, лог должен работать
          раньше GTK.
    \item \texttt{gtk\_init()}~--- инициализация GTK.
          Без неё нельзя создавать виджеты.
    \item \texttt{app\_data\_init()}~--- обнуление
          структуры \texttt{AppData}. Все поля — нули,
          \texttt{fd} = $-1$, буферы пусты.
    \item \texttt{load\_settings()}~--- чтение конфига.
          Заполняет \texttt{last\_device}, \texttt{auto\_reconnect}.
    \item \texttt{ui\_create()}~--- создание всех
          виджетов. Устанавливает \texttt{ui\_alive = TRUE}
          в конце.
    \item \texttt{scan\_serial\_ports()}~--- первое
          сканирование портов.
    \item \texttt{gtk\_widget\_show\_all()}~--- показ
          окна.
    \item \texttt{gtk\_main()}~--- главный цикл.
          Блокирует управление.
    \item \texttt{save\_settings()}~--- сохранение
          после выхода из главного цикла.
    \item \texttt{app\_data\_free()}~--- освобождение
          ресурсов.
    \item \texttt{return 0}~--- выход.
\end{enumerate}

\begin{keybox}
\textbf{Порядок инициализации~--- от \emph{общего} к \emph{частному}.}
Сначала~--- лог и GTK (инфраструктура). Потом~--- состояние
(\texttt{AppData}). Потом~--- конфиг. Потом~--- UI. Наконец~---
сканирование портов.
\end{keybox}

\begin{notebox}
\textbf{Порядок завершения~--- \emph{обратный} порядку
инициализации.} Сначала~--- сохраняем настройки. Потом~---
освобождаем ресурсы. Именно так: \emph{что} создали
\emph{последним}~--- освобождаем \emph{первым}.
\end{notebox}

\section{Инициализация: \texttt{app\_data\_init()}}
\label{sec:lifecycle-init}

Структура \texttt{AppData}~--- это всё состояние приложения.
Перед использованием её нужно обнулить.

\begin{listing}[htbp]
\begin{minted}[frame=none, bgcolor=gray!5]{c}
void app_data_init(AppData *app)
{
    memset(app, 0, sizeof(*app));
    app->fd = -1;
    app->auto_read_active = FALSE;
    app->is_connected = FALSE;
    plot_data_init(&app->plot_data);
    line_buffer_init(&app->line_buffer);
    g_mutex_init(&app->ring.lock);
}
\end{minted}
\caption{\texttt{app\_data\_init()}}
\label{lst:lifecycle-init}
\end{listing}

\subsection{Что происходит}

\textbf{Четыре шага:}

\begin{enumerate}
    \item \texttt{memset(app, 0, ...)}~--- обнуляем
          всю структуру. Все указатели становятся
          \texttt{NULL}, все числа~--- нулями.
    \item \texttt{app->fd = -1}~--- особый случай.
          По соглашению \texttt{fd < 0} означает
          <<порт не открыт>>. \texttt{memset} даст
          \texttt{fd = 0}~--- но это валидный файловый
          дескриптор! Поэтому явно ставим $-1$.
    \item \texttt{plot\_data\_init} и \texttt{line\_buffer\_init}~---
          инициализация вложенных структур.
    \item \texttt{g\_mutex\_init}~--- мьютекс для кольца
          кадров. Без него~--- гонка данных.
\end{enumerate}

\begin{warnbox}
\textbf{\texttt{fd = -1}~--- это \emph{соглашение}.} Если
забыть~--- \texttt{close\_serial\_port} может попытаться
закрыть \texttt{fd = 0} (что \emph{валидно}~--- но \emph{чужой}
дескриптор). Это приведёт к странным багам: <<порт закрыт,
а stdin тоже пропал>>.
\end{warnbox}

\begin{ideabox}
\textbf{Инициализация~--- это \emph{подготовка}, а не
\emph{настройка}.} Мы не \emph{задаём} \emph{рабочие}
значения~--- мы делаем структуру \emph{безопасной} для
\emph{последующего} использования. \emph{Реальные} значения
придут потом: из \texttt{load\_settings}, из
\texttt{open\_serial\_port}.
\end{ideabox}

\section{Загрузка настроек}
\label{sec:lifecycle-load}

\texttt{load\_settings()} вызывается сразу после
\texttt{app\_data\_init}. Она читает конфиг-файл
и заполняет поля \texttt{AppData}.

\begin{listing}[htbp]
\begin{minted}[frame=none, bgcolor=gray!5]{c}
void load_settings(AppData *app)
{
    char *path = config_path();
    GKeyFile *kf = g_key_file_new();

    if (!g_key_file_load_from_file(kf, path, G_KEY_FILE_NONE, NULL)) {
        app->auto_reconnect = TRUE;
        app->last_device[0] = '\0';
        return;
    }
    // ... чтение ключей ...
}
\end{minted}
\caption{\texttt{load\_settings()}~--- загрузка}
\label{lst:lifecycle-load}
\end{listing}

\subsection{Значения по умолчанию}

\textbf{Первый запуск}~--- конфига нет. Это нормально.
Приложение должно работать без него.

\begin{keybox}
\textbf{Три сценария:}
\begin{itemize}
    \item \textbf{Конфиг есть}~--- читаем.
    \item \textbf{Конфига нет}~--- значения по умолчанию.
    \item \textbf{Конфиг повреждён}~--- значения по умолчанию.
\end{itemize}
Все три должны приводить к \emph{работоспособному} приложению.
\end{keybox}

\section{Создание UI}
\label{sec:lifecycle-ui}

\texttt{ui\_create()}~--- самая большая функция.
Она создаёт все виджеты и связывает их
с колбэками.

\begin{listing}[htbp]
\begin{minted}[fontsize=\footnotesize, frame=none, bgcolor=gray!5]{c}
void ui_create(AppData *app)
{
    app->window = gtk_window_new(GTK_WINDOW_TOPLEVEL);
    // ... создание всех виджетов ...

    g_signal_connect(app->window, "delete-event",
                     G_CALLBACK(on_window_delete), app);
    g_signal_connect(app->window, "destroy",
                     G_CALLBACK(on_window_destroy), app);

    // ... создание панелей ...

    app->ui_alive = TRUE;
}
\end{minted}
\caption{\texttt{ui\_create()}}
\label{lst:lifecycle-ui}
\end{listing}

\subsection{Два сигнала на закрытие окна}

Обратите внимание: мы подключаем два разных
обработчика на закрытие окна.

\begin{itemize}
    \item \texttt{"delete-event"}~--- пользователь
          нажал крестик.
    \item \texttt{"destroy"}~--- окно уничтожается.
\end{itemize}

Это не дублирование~--- это разные моменты.
Между ними GTK уничтожает виджеты.

\begin{figure}[htbp]
\centering
\adjustbox{max width=\textwidth}{%
\begin{tikzpicture}[
    x=1cm, y=1cm,
    box/.style={
        draw, rounded corners=2mm, thick,
        minimum width=7cm, minimum height=0.8cm,
        align=center, font=\small
    },
    event/.style={box, fill=clrUIData!20},
    callback/.style={box, fill=clrLogicData!20},
    arr/.style={-{Latex[length=3mm]}, thick}
]

% ─── Координатная сетка (для отладки) ────────────
\draw[step=1cm, gray!20, very thin] (-1, -9) grid (10, 1);

\node[event]    (click)  at (4, 0)    {Пользователь жмёт крестик};
\node[event]    (delete) at (4, -1.2) {GTK посылает <<delete-event>>};
\node[callback] (cb1)    at (4, -2.4) {on\_window\_delete()};
\node[event]    (dw)     at (4, -3.6) {GTK уничтожает виджеты};
\node[event]    (destroy) at (4, -4.8) {GTK посылает <<destroy>>};
\node[callback] (cb2)    at (4, -6.0) {on\_window\_destroy()};
\node[event]    (end)    at (4, -7.2) {main() продолжается};

\draw[arr] (click)   -- (delete);
\draw[arr] (delete)  -- (cb1);
\draw[arr] (cb1)     -- (dw);
\draw[arr] (dw)      -- (destroy);
\draw[arr] (destroy) -- (cb2);
\draw[arr] (cb2)     -- (end);

% Метки осей (для отладки)
\foreach \x in {-1,...,10} {
    \node[font=\tiny, gray] at (\x, -9.3) {\x};
}
\foreach \y in {-9,...,1} {
    \node[font=\tiny, gray] at (-1.3, \y) {\y};
}

\end{tikzpicture}
}% end adjustbox
\caption{Жизненный цикл окна при закрытии.}
\label{fig:lifecycle-window}
\end{figure}

\section{Завершение: \texttt{app\_data\_free()}}
\label{sec:lifecycle-free}

\texttt{app\_data\_free()}~--- это зеркало
\texttt{app\_data\_init}. Она освобождает всё,
что было создано.

\begin{listing}[htbp]
\begin{minted}[frame=none, bgcolor=gray!5]{c}
void app_data_free(AppData *app)
{
    stop_auto_reconnect(app);

    if (app->redraw_timer_id) {
        g_source_remove(app->redraw_timer_id);
        app->redraw_timer_id = 0;
    }

    if (app->osc_thread) {
        app->thread_running = false;
        g_thread_join(app->osc_thread);
        app->osc_thread = NULL;
    }

    app->ui_alive = FALSE;
    close_serial_port(app);

    g_mutex_clear(&app->ring.lock);
}
\end{minted}
\caption{\texttt{app\_data\_free()}}
\label{lst:lifecycle-free}
\end{listing}

\subsection{Порядок остановки}

\textbf{Порядок}~--- критичен:

\begin{enumerate}
    \item \texttt{stop\_auto\_reconnect()}~--- сначала
          останавливаем таймер реконнекта. Иначе
          он может сработать после того, как
          мы уничтожим \texttt{AppData}.
    \item \texttt{g\_source\_remove(redraw\_timer\_id)}~---
          останавливаем таймер перерисовки.
          Тот же принцип.
    \item \texttt{g\_thread\_join()}~--- останавливаем
          поток чтения. Пока он работает~---
          нельзя освобождать кольцо.
    \item \texttt{app->ui\_alive = FALSE}~--- помечаем UI
          мёртвым. Это защита для \texttt{close\_serial\_port}.
    \item \texttt{close\_serial\_port()}~--- закрываем
          порт последним.
    \item \texttt{g\_mutex\_clear}~--- освобождаем мьютекс.
\end{enumerate}

\begin{warnbox}
\textbf{Порядок остановки~--- \emph{обратный} порядку
создания.} Что \emph{создали} \emph{последним}~--- останавливаем
\emph{первым}. Что \emph{создали} \emph{первым}~---
останавливаем \emph{последним}.
\end{warnbox}

\section{Полный жизненный цикл}
\label{sec:lifecycle-full}

Соберём всё вместе: запуск, работа,
завершение.

\begin{figure}[htbp]
\centering
\adjustbox{max width=\textwidth}{%
\begin{tikzpicture}[
    x=1cm, y=1cm,
    box/.style={
        draw, rounded corners=1.5mm, thick,
        minimum width=4.5cm, minimum height=0.8cm,
        align=center, font=\footnotesize\ttfamily,
        fill=white
    },
    phase/.style={
        draw, rounded corners=3mm, thick,
        minimum width=3cm, minimum height=0.8cm,
        align=center, font=\small\bfseries,
        fill=clrIdea!10
    },
    arr/.style={-{Latex[length=3mm]}, thick}
]

% ─── Координатная сетка (для отладки) ────────────
\draw[step=1cm, gray!20, very thin] (-1, -19) grid (10, 1);

% Запуск
\node[phase] (start) at (4, 0) {ЗАПУСК};
\node[box, below=of start] (init) {app\_data\_init()};
\node[box, below=of init] (load) {load\_settings()};
\node[box, below=of load] (ui) {ui\_create()};
\node[box, below=of ui] (scan) {scan\_serial\_ports()};
\node[box, below=of scan] (show) {gtk\_widget\_show\_all()};

% Работа
\node[phase, below=1cm of show] (run) {РАБОТА};
\node[box, below=of run] (main) {gtk\_main()};

% Завершение
\node[phase, below=1cm of main] (end) {ЗАВЕРШЕНИЕ};
\node[box, below=of end] (save) {save\_settings()};
\node[box, below=of save] (free) {app\_data\_free()};

\draw[arr] (start) -- (init);
\draw[arr] (init) -- (load);
\draw[arr] (load) -- (ui);
\draw[arr] (ui) -- (scan);
\draw[arr] (scan) -- (show);
\draw[arr] (show) -- (run);
\draw[arr] (run) -- (main);
\draw[arr] (main) -- (end);
\draw[arr] (end) -- (save);
\draw[arr] (save) -- (free);

% Метки осей (для отладки)
\foreach \x in {-1,...,10} {
    \node[font=\tiny, gray] at (\x, -19.3) {\x};
}
\foreach \y in {-19,...,1} {
    \node[font=\tiny, gray] at (-1.3, \y) {\y};
}

\end{tikzpicture}
}% end adjustbox
\caption{Полный жизненный цикл приложения.}
\label{fig:lifecycle-full}
\end{figure}

\subsection{Три фазы}

\begin{enumerate}
    \item \textbf{Запуск}~--- от \texttt{main()} до
          \texttt{gtk\_widget\_show\_all}. \emph{Создаём}
          всё: состояние, UI, таймеры.
    \item \textbf{Работа}~--- \texttt{gtk\_main()}.
          Обработка событий: клики, таймеры, данные из порта.
          Пока пользователь не закроет окно~---
          \texttt{gtk\_main} \emph{не вернёт} управление.
    \item \textbf{Завершение}~--- от \texttt{save\_settings}
          до \texttt{return 0}. \emph{Освобождаем}
          всё в \emph{обратном} порядке.
\end{enumerate}

\begin{keybox}
\textbf{\texttt{gtk\_main()}~--- это <<сердце>> приложения.}
Он \emph{блокирует} управление. Всё, что происходит
\emph{во время} работы, происходит \emph{внутри}
\texttt{gtk\_main}~--- через \emph{колбэки}.
\end{keybox}

\section{Уроки главы}
\label{sec:lifecycle-lessons}

\begin{enumerate}
    \item \textbf{Порядок инициализации~--- от общего
          к частному.} Лог, GTK, состояние, конфиг, UI.
    \item \textbf{Порядок завершения~--- обратный.}
          Конфиг, состояние, GTK.
    \item \textbf{\texttt{app\_data\_init}~--- не <<просто
          memset>>.} \texttt{fd = -1}~--- \emph{соглашение}.
    \item \textbf{Приложение должно работать без конфига.}
          Первый запуск~--- конфига нет. Это нормально.
    \item \textbf{Два сигнала на закрытие окна~--- не
          дублирование.} \texttt{delete-event}~--- \emph{до}
          уничтожения виджетов, \texttt{destroy}~--- \emph{после}.
    \item \textbf{Таймеры~--- первыми.} \texttt{app\_data\_free}
          сначала останавливает таймеры, потом~--- порт.
    \item \textbf{\texttt{gtk\_main}~--- сердце приложения.}
          Всё происходит внутри него.
    \item \textbf{Потоки~--- до освобождения памяти.}
          \texttt{g\_thread\_join} перед \texttt{g\_mutex\_clear}.
\end{enumerate}

\section{Чек-лист грабель}
\label{sec:lifecycle-pitfalls}

\begin{table}[htbp]
\centering
\caption{Грабли жизненного цикла}
\label{tab:lifecycle-pitfalls}
\small
\begin{tabularx}{\textwidth}{@{}c L L@{}}
\toprule
\textbf{№} & \textbf{Грабля} & \textbf{Решение} \\
\midrule
1 & \texttt{debug\_log\_init} после \texttt{gtk\_init} & До~--- иначе ошибки GTK потеряются \\
2 & \texttt{fd = 0} по умолчанию & Явно ставить $-1$ в \texttt{app\_data\_init} \\
3 & Приложение падает без конфига & Значения по умолчанию \\
4 & Закрытие порта в \texttt{destroy} & В \texttt{delete-event} \\
5 & Таймеры работают после \texttt{app\_data\_free} & Остановить \emph{до} освобождения \\
6 & \texttt{save\_settings} до \texttt{gtk\_main} & После~--- иначе нечего сохранять \\
7 & \texttt{app\_data\_free} без \texttt{close\_serial\_port} & Обязательно закрыть порт \\
8 & Поток работает после \texttt{g\_mutex\_clear} & \texttt{g\_thread\_join} \emph{до} \\
\bottomrule
\end{tabularx}
\end{table}

\section{Упражнения}
\label{sec:lifecycle-exercises}

\begin{exercisebox}
\textbf{Упражнение 1.} Что произойдёт, если поменять местами
\texttt{gtk\_init} и \texttt{debug\_log\_init} в \texttt{main}?
Приведите сценарий, в котором это приведёт к потере
информации.
\end{exercisebox}

\begin{exercisebox}
\textbf{Упражнение 2.} Добавьте в \texttt{main} обработку
аргументов командной строки: флаг \texttt{--port /dev/ttyUSB0}
для автоматического подключения. Как это повлияет на
\texttt{load\_settings}?
\end{exercisebox}

\needspace{5\baselineskip}
\begin{exercisebox}
\textbf{Упражнение 3.} Реализуйте корректную обработку
\texttt{SIGINT} (Ctrl+C из терминала). Что должно произойти
до выхода? Как связать сигнал с \texttt{gtk\_main\_quit}?
\end{exercisebox}

\begin{exercisebox}
\textbf{Упражнение 4.} Что произойдёт, если \texttt{app\_data\_free}
вызывается \emph{до} \texttt{gtk\_main}? Приведите пример,
когда это \emph{допустимо}.
\end{exercisebox}

\begin{exercisebox}
\textbf{Упражнение 5.} Нарисуйте (на бумаге или в TikZ) полный
жизненный цикл \emph{вашего} приложения. Какие модули в каком
порядке инициализируются и освобождаются?
\end{exercisebox}


